\section{The same native function fields from an actual nonempty isomorphism open}
\label{sec:generic-open-function-field}

\begin{theorem}[Actual dominance and original stalk isomorphisms]
\label{sc0-generic-open-field:original-dominance-stalks}
\lean{EnvelopeProperty.GenericOpenFunctionField.isDominant_of_isIso_restrict,EnvelopeProperty.GenericOpenFunctionField.isIso_stalkMap_of_isIso_restrict}
\leanok
Let $g:Y\to X$ be the original scheme map, and let $U\subset X$ be
an original nonempty open over which its actual restriction is an
isomorphism. If $X$ is integral, the SAME original $g$ is dominant.
For every original $y$ over $U$, its ENTIRE original native stalk map
is an isomorphism. This latter assertion needs neither integrality
nor nonemptiness of $U$ as an additional premise.
\end{theorem}
\begin{proof}
For every original point of $U$, the actual inverse restricted
isomorphism provides a point of the original inverse-image open.
The genuine original restriction triangle shows that its image under
the SAME original map is the original point. Thus $U$ is contained
in the actual range of $g$. A nonempty open of an integral scheme
is dense, so the entire original range is dense.

For a point of the original preimage, the actual native restriction
of the original map is an isomorphism and hence has an isomorphism
on its original stalks. The native restriction arrow isomorphism
identifies that ENTIRE map with the SAME original stalk map, using
the actual open-inclusion stalk comparisons. Transport of the
isomorphism property through that genuine arrow isomorphism proves
the assertion, with no supplied local-map comparison certificate.
\end{proof}

\begin{theorem}[The same entire original generic-stalk field isomorphism]
\label{sc0-generic-open-field:whole-original-fields}
\lean{EnvelopeProperty.GenericOpenFunctionField.isIso_functionFieldMap_of_isIso_restrict,EnvelopeProperty.GenericOpenFunctionField.functionFieldIso,EnvelopeProperty.GenericOpenFunctionField.functionFieldIso_hom,EnvelopeProperty.GenericOpenFunctionField.functionFieldIso_section}
\leanok
If $Y$ is also integral, the SAME actual native dominant generic-stalk
map is an isomorphism between the original function fields. The
constructed function-field isomorphism has that ENTIRE original map
as its forward morphism. On EVERY original nonempty open whose
preimage is nonempty, it sends the actual generic germ of EVERY
original coefficient section to the actual generic germ of its SAME
original pullback section.
\end{theorem}
\begin{proof}
The actual restricted isomorphism and the nonempty original $U$
give a nonempty original preimage open. The original generic point
of the integral $Y$ lies in this open. The preceding original stalk
isomorphism theorem therefore applies at that SAME original generic
point. Actual dominance maps the original generic points to one
another. The native function-field map is exactly this actual stalk
map preceded by the original equality-induced stalk specialization.
That specialization is an isomorphism: substituting the actual point
equality reduces it to the original identity. Their genuine composite
is therefore an isomorphism.

Take the actual categorical isomorphism of this SAME map, whose
forward morphism is definitionally the original generic-stalk map.
The previously proved actual whole generic-stalk coefficient square
then gives the EVERY-section formula, with the original coefficient
and the SAME original scheme pullback retained on both sides. No
abstract field equivalence, dominance premise, field-map matching
or coefficient certificate is supplied.
\uses{sc0-generic-open-field:original-dominance-stalks}
\end{proof}

\begin{theorem}[The same originally specified affine fraction field]
\label{sc0-generic-open-field:specified-original-fractions}
\lean{EnvelopeProperty.GenericOpenFunctionField.specFractionFieldEquiv,EnvelopeProperty.GenericOpenFunctionField.specFractionFieldEquiv_pullback_section,EnvelopeProperty.GenericOpenFunctionField.specFractionFieldEquiv_coefficients}
\leanok
If the original base is $\operatorname{Spec}A$ for a domain $A$, and
$F$ is its originally specified fraction field, the actual target
function field has a genuine ring equivalence to that SAME $F$.
Every original pulled coefficient section has the value obtained
from its original base generic germ by the actual fraction-field
equivalence. In particular, EVERY original $a\in A$ has exactly its
prescribed value under the original map $A\to F$.
\end{theorem}
\begin{proof}
Invert the SAME original generic-stalk isomorphism from the preceding
theorem and compose with the genuine equivalence of the native
affine generic-stalk fraction field and the specified original $F$.
Apply the true whole inverse identity to the preceding EVERY-section
comparison. This gives the desired coefficient formula for every
actual section. For an original $a\in A$, the original native generic
germ of its actual spectrum section is the native affine algebra map
of $a$. The actual fraction-field algebra equivalence preserves that
SAME original coefficient, giving precisely its original value in
$F$. No fraction-field or coefficient-matching certificate is added.
\uses{sc0-generic-open-field:whole-original-fields}
\end{proof}

These are ordinary geometric prerequisites for comparing the actual
BP target field with the original affine fraction field. They do not
assert regularity of arbitrary pushed relative canonical coefficients,
any full rounded canonical MODULE pushforward formula, or vanishing.
